\chapter{实验设计}

\section{什么才算本报告中的实验}

本报告只纳入能回答以下至少一个问题的实验：

\begin{enumerate}
  \item 正式模型是否真实训练、是否持续降低示教动作误差？
  \item 团队探索了多少数据、相机、资源和训练方案，哪些运行真正走得足够远？
  \item 数据是否覆盖无盖、方向、翻转、回位和自由旋转瓶盖等关键条件？
  \item 模型是否真实利用三路视野？
  \item 强化学习研究是否形成可训练、可保存、可离线验证的闭环？
\end{enumerate}

没有目的、没有对照、只因为“文件存在”的测试不进入结果章节。训练进程标记结束不等于模型有效；单条示教关键帧不等于机器人成功；训练损失下降不等于任务成功率。

\section{训练试验谱系}

对 2026 年 4 月 1 日至 5 月 31 日的训练平台记录按瓶子分拣、拧盖、冲洗和插入等明确关键词筛选，得到 41 次运行尝试、33 个不同运行名称和 14 组配置。其中基础瓶子分拣方向 23 次，冲洗/插入探索 18 次。\ev{E13}

\figfull[.96\textwidth]{experiment_funnel.pdf}{训练试验漏斗。该图衡量工程迭代量和运行稳定性，不衡量模型真机能力。}

\begin{table}[H]
\centering
\caption{训练试验矩阵}
\begin{tabularx}{\textwidth}{>{\bfseries}p{34mm}p{23mm}p{29mm}p{28mm}X}
\toprule
方向 & 尝试数 & 不同配置 & 达到1万步 & 目的与限制\\
\midrule
瓶子分拣基础模型 & 23 & 10 & 9 & 探索批量规模、全参数/轻量训练、相机与数据组合；仅一个随机种子\\
冲洗/插入探索 & 18 & 4 & 3 & 探索冲洗和后续插入任务；目标不同，损失不可与分拣直接比较\\
合计 & 41 & 14 & 12 & 33 次中断、5 次失败、3 次标记结束；“结束”不等于成功\\
\bottomrule
\end{tabularx}
\end{table}

运行批量规模覆盖 4、8、32、128、256、512，说明大量工作用于解决显存、数据读取、分布式训练和稳定性问题。所有可识别运行只使用随机种子 42，尚不能报告多种子均值或置信区间。

\section{正式模型训练实验}

正式训练实验使用 25 个唯一数据仓库、29 个采样条目、批量规模 256 和 4 张 H200，留下 6,059 条历史记录。评估项目包括训练损失、梯度范数、参数范数、数据等待时间和单步时间；没有验证损失、测试集指标或真机任务结果。\ev{E08}

\begin{table}[H]
\centering
\caption{正式训练实验的可用指标与不可用指标}
\begin{tabularx}{\textwidth}{>{\bfseries}p{38mm}p{24mm}X}
\toprule
指标 & 是否记录 & 能回答的问题\\
\midrule
训练动作生成损失 & 是 & 模型是否越来越接近示教动作\\
梯度/参数范数 & 是 & 训练是否数值稳定\\
训练时长与步数 & 是 & 训练工作量和保存点位置\\
验证集损失 & 否 & 不能判断训练/验证差距\\
测试集成功率 & 否 & 不能判断未见场景泛化\\
真机阶段完成率 & 否 & 不能判断失败发生在哪个阶段\\
正式吞吐量 & 否 & 现场约 2 瓶/分钟只能作估计\\
\bottomrule
\end{tabularx}
\end{table}

\section{现场展示观测}

现场展示中观察到完整取瓶—开盖—瓶/盖分箱循环；观察者估计连续工作时间超过 1 小时、平均约 2 瓶/分钟，并认为系统对桌面随机位置和一定瓶子变化有初步适应能力。\ev{E10}

本次资料没有完整录像、计数表、逐次测试记录、固定初始条件和失败时间戳。因此现场展示只按“观测估计”报告，不计算成功率、标准差或置信区间；也不把“超过一小时”与注意力清单的时间跨度相互替代。

\section{强化学习探索实验}

强化学习研究使用另一组“冲洗瓶子”数据：835 条轨迹、354,835 个原始帧，经裁剪后 299,347 帧，形成 238,816 个训练片段；人工终止标签中 137 条为正、698 条为零。\ev{E15} 这些标签来自人工编辑，不是自动物理成功检测，也不是基础瓶子分拣模型的真机成功率。

连续研究轮次 28--32 使用固定的 11 条离线验证轨迹，没有独立测试集。最后一轮动作优化/价值判断保存模型包含 61,541,926 个参数，优化器状态和目标网络均存在。\ev{E16}

\begin{table}[H]
\centering
\caption{强化学习探索的证据成熟度}
\begin{tabularx}{\textwidth}{>{\bfseries}p{31mm}p{24mm}X}
\toprule
环节 & 状态 & 说明\\
\midrule
真实数据与人工奖励 & 已完成 & 835 条轨迹；奖励质量仍依赖人工标注\\
训练数据生成 & 已完成 & 238,816 个可训练片段\\
模型训练与保存 & 已完成 & 连续 5 个轮次、保存参数可读取\\
固定离线验证 & 已完成 & 11 条；规模小且无独立测试\\
同条件基础模型对照 & 未完成 & 无法判断强化学习是否真正优于基础模型\\
真机增益验证 & 未完成 & 不能写成已提高任务成功率\\
\bottomrule
\end{tabularx}
\end{table}

\begin{table}[H]
\centering
\caption{强化学习结果进入本报告的判定}
\begin{tabularx}{\textwidth}{>{\bfseries}p{38mm}p{24mm}X}
\toprule
候选结论 & 是否纳入 & 理由\\
\midrule
已建立真实数据到训练保存的闭环 & 纳入 & 数据、模型和日志可交叉核验\\
固定离线验证动作误差下降 & 纳入 & 同一验证清单、连续轮次可追溯\\
人工正标签占137/835 & 仅作数据说明 & 不是自动物理成功率，也不是基础分拣任务结果\\
强化学习已经提高真机成功率 & 不纳入 & 缺同条件基础模型和真机对照\\
单条正标签轨迹证明方法有效 & 不纳入 & 单案例不能代表总体能力\\
\bottomrule
\end{tabularx}
\end{table}
